\scanpage{71}\mspage{F.}
\margin{\uline{Pic birationnel}}
% slanted
\uline{Théorème} Soient $f:X\to Y$ un morphisme \inserted{birationnel} dominant
de préschémas intègres normaux, \deleted{et} propres sur $k$,
avec $Y$ \deleted{loc. \uline{factoriel}}\inserted{\uline{régulier}}. Alors le morphisme
\[
\underline{\mathrm{Pic}}^{\tau}_Y\longrightarrow\underline{\mathrm{Pic}}^{\tau}_X
\]
est un \uline{isomorphisme}.
\margin{$X$ normal ``sur $k$'' ?}
% slanted
\margin{\unsure{faut-il|peut-on} $=$ \deleted{\unsure{supposer}} $Y$ \deleted{régulier ??} \deleted{\unsure{non}} \unsure{simple|sing.} sur $k$ ?}
% slanted
\margin{N.B. factoriel ne suffit pas $=$ du pt de \unsure{vue} \unsure{ensembliste|ensemblips}, cf cônes ---}
% slanted

C'est un \uline{monomorphisme} grâce au fait que
$f_*(\underline{O}_X)=\underline{O}_Y$ (universellement sur $k$ !), donc une
immersion fermée puisque un hom de groupes
-- ou \unsure{parce que} propre. Montrons que c'est
\uline{surjectif}. Soit $U$ l'ouvert de $Y$ où $f^{-1}$ est
défini, comme $Y$ est factoriel, on a
$\mathrm{Pic}(X)\overset{\sim}{\longrightarrow}\mathrm{Pic}(U)$, d'où, utilisant
\diagram{\begin{tikzcd}[ampersand replacement=\&,column sep=small,row sep=small]
{}\&X\arrow[d,"f"]\\
U\arrow[ur,"g"]\arrow[r,"i"]\&Y
\end{tikzcd}}
, le diagramme commutatif
\diagram{\begin{tikzcd}[ampersand replacement=\&,column sep=small,row sep=small]
{}\&\mathrm{Pic}(X)\arrow[dl,leftrightarrow,"g^*"]\\
\mathrm{Pic}(U)\&\mathrm{Pic}(Y)\arrow[l,"i^*","\sim"']\arrow[u,"f^*"']
\end{tikzcd}}
qui montre que $\mathrm{Pic}(Y)\to\mathrm{Pic}(X)$ est un hom
sur un facteur direct, l'image étant \unsure{formé|formée} des
$\xi$ telles que $\deleted{\unsure{g^*(\xi)=0.|f^*(\xi)=0}}\ \xi=f^*i^{*-1}g^*(\xi)$. Or
$\xi$ est \deleted{\unsure{num}} $\tau$-équivalente à $0$, \inserted{donc $\exists\,\unsure{n|m}\xi$ est alg$^{\mathrm t}$ équiv. à $0$,} \deleted{\unsure{son}} $g^*(\xi)$
\deleted{est} \deleted{$\tau$-équivalent}\inserted{alg. équiv.} à $0$, donc $i^{*-1}g^*(\xi)$ \deleted{\unsure{est}} aussi,
donc \deleted{\unsure{$n$}} $f^*i^{*-1}g^*(\xi)$ \inserted{$=n f^*(i^{*-1}g^*(\xi))$} \unsure{égal$^{\mathrm t}$}. D'autre part, on voit que
$\xi-f^*i^{*-1}\deleted{g^*(\xi)}$ \deleted{\unsure{est diviseur représenté il}} est représenté par un
\margin{\unsure{Revenir au dessin|Revenir au discours|Raisonner en termes} des cycles !}
% slanted
\margin{$Y$ \unsure{régulier|séparation} ?? ou \unsure{simple|tjrs simple} sur $k$}
% slanted
% archive sheet number: 70
